\section{Discussion and Production Implications}
\label{sec:discussion}

BanditGPT bridges the gap between theoretical contextual bandit algorithms and the practical realities of deploying multi-LLM routers in production environments. We highlight several key implications for engineering teams and address current limitations.

\subsection{Practical Impact}

\paragraph{Interpretable Interface and SLA Compliance.} 
The standard approach to cost-aware routing uses a static penalty $\lambda$, which requires translating a dimensionless weight into a dollar outcome via offline simulation. This mapping breaks whenever the portfolio changes. BanditGPT's Primal-Dual \texttt{BudgetPacer} accepts a dollar-denominated target (e.g., $B = \$0.002$/request) directly. This target is portable across portfolios because the normalized dual update ($\bar{c}_t / B$) adapts the penalty $\lambda_t$ to the realized cost scale automatically. 

Furthermore, the pacer's precise budget utilization ($0.97\times$ to $1.02\times$) allows it to underpin strict Service Level Agreements (SLAs) for cost (e.g., ``average spend will not exceed $1.10 \times B$ over any rolling window''). Static penalties offer no such guarantee, as realized costs fluctuate arbitrarily if prompt distributions drift.

\paragraph{Eliminating Offline Recalibration.}
In a deployment where model pricing changes frequently (as is common with frontier APIs), a static penalty becomes stale, requiring costly offline retraining to discover the new optimal $\lambda$. The \texttt{BudgetPacer} automatically detects price drops via its online cost EMA and dynamically reallocates freed budget to higher-quality models. This autonomous optimization eliminates the need for manual intervention, maximizing application quality while strictly obeying financial constraints.

\paragraph{Cost Savings at Scale.}
Consider a deployment routing 100,000 requests per day. Moving from an unconstrained policy (\$0.0094/req, quality $0.921$) to the pacer at $B{=}\$0.0019$/req saves \$760/day (over \$277,000/year) while maintaining quality of $0.908$---a 1.4\% quality reduction for an 81\% cost reduction.

\paragraph{Portfolio and Encoder Portability.}
The experiments in this paper evaluate a fixed configuration (MiniLM-L6 encoder, $K{=}3$ portfolio, shipped warmup priors), but production teams operate diverse model portfolios that evolve over time---models are added when new APIs launch, removed when deprecated, and swapped when pricing changes.
A router that is hard-wired to a fixed encoder or model set forces a full retrain whenever the portfolio changes, negating the benefit of online adaptation.

BanditGPT decouples each component of the pipeline so that every part can be replaced independently:
\begin{itemize}[nosep,leftmargin=*]
  \item \textbf{Encoder:} Teams with domain-specific embeddings (e.g., code, biomedical, multilingual) can supply any \texttt{Callable[[str], np.ndarray]} as a custom encoder, or pass pre-computed vectors directly, without depending on \texttt{sentence-transformers}.
  \item \textbf{PCA and warmup priors:} A two-function calibration API (\texttt{train\_pca}, \texttt{generate\_warmup\_priors}) produces compatible artifacts for any encoder--portfolio combination from a labelled reward dataset.  This lets a team onboard a new portfolio in minutes rather than days---encode representative prompts, collect reward labels, and generate priors---without modifying the routing algorithm.
  \item \textbf{Model portfolio:} Models can be added or removed at runtime via \texttt{add\_model()} and \texttt{remove\_model()} without restarting the service.  New arms start from the isotropic prior and begin learning immediately; removed arms are garbage-collected from the bandit state.
\end{itemize}
This portability matters because the most common deployment failure mode for adaptive routers is not algorithmic---it is operational friction: the effort required to adapt the system when the model landscape changes.  By making every artifact regenerable and every component swappable, BanditGPT ensures that the routing algorithm remains the stable core while the surrounding model ecosystem evolves freely.

\subsection{Limitations and Deployment Considerations}

\paragraph{Warm-Start Dependence at Tight Budgets.}
At tight budgets, the router selects expensive arms rarely during Phase~1 ($<1\%$ Gemini at $B{=}\$2.3{\times}10^{-4}$), so the bandit's quality estimate for those arms relies almost entirely on the warmup prior rather than online observations.
In our experiments this is not limiting---the warm prior is sufficiently accurate and the primal-dual mechanism rapidly redirects traffic once the price drops (Section~\ref{sec:results})---but the dependence on prior quality is stronger than at moderate or loose budgets.
If the warmup prior were poorly calibrated for a newly added model (e.g., no historical data), a tight budget could delay adoption even after a price drop.
Future work could investigate targeted exploration strategies (e.g., small $\epsilon$-greedy budget reserves) that periodically evaluate under-explored arms to maintain calibrated quality estimates independent of the prior.

\paragraph{Static Unit-Cost Normalization.}
The normalized cost $\tilde{c}_a$ in Equation~\ref{eq:budget_ucb} and the hard ceiling filter both use a static blended rate $(\text{input\_cost} + \text{output\_cost})/2$ rather than weighting by the actual input/output token counts of each request.
For workloads with highly asymmetric token ratios (e.g., retrieval-augmented generation with $\gg$10$\times$ more input than output tokens), a token-weighted cost would give a more precise per-request estimate.
In practice, this does not affect our experimental conclusions for three reasons: (i)~the Pareto frontiers are evaluated using \emph{exact per-request costs} from the offline dataset, not the router's blended estimate; (ii)~the log-scale normalization preserves the cost \emph{ranking} of models across any reasonable token ratio; and (iii)~the bandit's learned quality model is independent of the cost penalty---$\tilde{c}_a$ biases arm selection but does not distort the reward posterior.
A production deployment serving heterogeneous workloads (e.g., both summarization and RAG) could benefit from passing the expected token ratio at routing time; the open-source release already accepts \texttt{input\_tokens}/\texttt{output\_tokens} parameters and uses them for post-hoc cost logging.

\paragraph{Obtaining Discriminative Rewards in Practice.}
Our experiments use an automated LLM judge that produces continuous, multi-dimensional scores---an ideal signal for $K{\geq}3$ routing, where the bandit must distinguish models whose average quality differs by only a few percentage points (e.g., $0.908$ vs.\ $0.921$ in our portfolio).
Binary feedback (thumbs-up/down) is insufficient in this regime: as shown in Experiment~1, a natural binarisation threshold ($\tau{=}0.85$) causes the two stronger models to tie on $85.8\%$ of prompts, collapsing the three-arm problem to two and eliminating the router's ability to find quality within a budget.
This finding is consistent with the broader literature: production routers serving $K{>}2$ portfolios universally rely on continuous quality prediction~\cite{wang2025mixllm,bhatti2026proteus}, while binary feedback is restricted to $K{=}2$ triage systems~\cite{ong2024routellm,wang2025barp}.

In customer-facing deployments where a dedicated LLM judge is too costly or too slow, four practical strategies can produce the continuous reward signal that multi-model routing requires:
\begin{enumerate}[nosep,leftmargin=*]
  \item \textbf{Asynchronous LLM judge (sampled).}  Rather than evaluating every response, teams can judge a random sample (e.g., 5--10\%) asynchronously and use the resulting scores as stochastic rewards.  LinUCB converges under partial feedback; the effective sample rate simply slows learning proportionally---acceptable when warmup priors already provide a strong initialisation.
  \item \textbf{Multi-criteria human ratings.}  Replacing a binary thumbs-up/down with a short Likert scale (e.g., 1--5 stars) or a two-question rubric (``Was this helpful?'' + ``Was this complete?'') produces a $[0,1]$ composite that preserves inter-model gaps.  Even a three-point scale (poor/acceptable/excellent) retains sufficient discrimination for $K{=}3$--$5$ portfolios.
  \item \textbf{Implicit behavioural composites.}  Many applications already log signals that correlate with response quality: whether the user copied the output, the time spent reading it, follow-up question rate, code acceptance or edit distance, task completion flags, or retry/regeneration rate.  Individually each signal is noisy and near-binary, but combining $m$ such indicators into a normalised composite $r_t = \frac{1}{m}\sum_i s_i \in [0,1]$ yields a continuous reward whose variance decreases as $1/m$.  This approach requires no explicit rating UI and is compatible with high-volume, low-touch applications.
  \item \textbf{Windowed binary aggregation.}  If only binary feedback is available, the per-arm success proportion over a rolling window (e.g., last 200 observations per arm) approximates a continuous expected reward.  However, this strategy requires substantially more traffic per arm before inter-model gaps become statistically detectable, making it viable only for high-volume deployments.
\end{enumerate}
BanditGPT's reward interface accepts any scalar $r_t \in [0,1]$, so all four strategies---and hybrid combinations---are drop-in compatible without modifying the routing algorithm.

\paragraph{Delayed Feedback and State Persistence.}
Our experiments assume immediate feedback, but three common production scenarios create a gap between routing and reward observation.
First, \emph{automated judge evaluation} (as used in this paper) is itself an LLM call that adds latency and cost; production systems typically batch these evaluations asynchronously rather than blocking the response path.
Second, \emph{human feedback} (thumbs-up/down, post-edit acceptance, task-completion signals) is inherently delayed---users do not rate responses immediately, and downstream quality signals such as conversion or error rates may take hours or days to materialize.
Third, in microservice architectures the \emph{routing service and feedback-ingestion service may be separate processes} or even separate machines, so the context vector $x_t$ used at route time is no longer in memory when the reward $r_t$ arrives.

LinUCB's update rule requires the \emph{exact} context vector that was used at selection time; re-encoding the prompt is not reliable because the embedding model version, PCA artifact, or whitening scales may have changed between routing and feedback ingestion.
The context vector must therefore be persisted at route time and retrieved by request ID when feedback arrives.
Without such persistence, a process restart (deployment, crash, autoscaling event) also discards all learned state, forcing the router to cold-start from priors and re-learn from scratch.
Existing bandit-based routers~\cite{panda2025pilot,wang2025barp,li2025llmbandit} maintain state only in process memory, making them incompatible with all three scenarios above.

BanditGPT addresses this with a pluggable \texttt{ContextStore} abstraction backed by two built-in implementations:
\begin{itemize}[nosep,leftmargin=*]
  \item \texttt{EphemeralContextStore}: A bounded in-memory deque for low-latency deployments where evaluation is synchronous (e.g., automated metrics computed inline).
  \item \texttt{SqliteContextStore}: A disk-persisted store using WAL-mode SQLite with TTL-based expiration, supporting millions of entries, concurrent reads/writes ($>$10K writes/sec), and week-scale retention for RLHF workflows---with zero external dependencies.
\end{itemize}
A separate \texttt{CheckpointManager} periodically snapshots the learned bandit state ($A_a$, $A_a^{-1}$, $b_a$, $\lambda_t$) via atomic writes, enabling crash recovery with automatic registry-drift handling when the model portfolio changes between restarts.
The storage layer is extensible: implementing the \texttt{ContextStore} interface for Redis or S3 requires overriding three methods without modifying router logic.

\paragraph{Embedding Generation Costs.}
Contextual bandits require a feature vector $x_t$ for every prompt. Generating this vector using a dense embedding model (e.g., SentenceTransformers) incurs a small latency and compute cost. While negligible compared to the inference cost of frontier LLMs, this overhead must be factored into the overall budget $B$ for ultra-low-cost, high-throughput deployments. Utilizing smaller, faster embedding models or distilling features directly from the prompt text are practical mitigations.

\section{Conclusion}

We introduced BanditGPT, an open-source framework for budget-constrained, non-stationary LLM routing.  By framing the problem as a Primal-Dual Contextual Bandit with Knapsacks and integrating geometric forgetting, we provide a robust solution to the dynamic challenges of production deployments.  Our system demonstrates precise budget compliance ($0.97$--$1.02\times$ target), eliminates the need for manual hyperparameter tuning of static cost penalties, and autonomously adapts to silent model updates and mid-cycle price drops.

\paragraph{Why this matters for practitioners.}
The business case for adaptive, multi-model routing rests on three converging operational realities.
First, \emph{budgets are binding}: as LLM spend grows from experimental line items to material operating costs, finance and engineering teams require per-request SLAs rather than after-the-fact reconciliation.  A router that can draw from $K{\geq}3$ models at distinct price points has more degrees of freedom to find quality \emph{within} a budget---it can route easy prompts to a mid-tier model that matches frontier quality at a fraction of the cost, stretching the budget to afford frontier calls where they matter most.  With only $K{=}2$, the router's sole lever is the fraction of traffic sent to the expensive model.
Second, \emph{model quality is non-stationary}: providers silently update weights, deprecate endpoints, and release improved alternatives on a monthly cadence.  Every such event can degrade the current best arm; a larger portfolio increases the chance that an alternative maintains quality at a comparable cost, giving the adaptive router more options to recover without violating its budget constraint.
Third, \emph{costs shift frequently}: when a provider cuts prices, a richer portfolio creates more opportunities for the router to reallocate freed budget toward quality improvements.  Our experiments show that BanditGPT's BudgetPacer automatically detects such opportunities and re-concentrates traffic on the newly cost-effective arm within hundreds of requests (Section~\ref{sec:results}).

BanditGPT is designed to operate at this intersection of budget constraints, reward shifts, and cost shifts across multi-model portfolios.  Beyond the algorithmic contributions, the library provides the production scaffolding these scenarios demand---pluggable encoders and PCA artifacts for domain-specific embeddings, persistent context stores for delayed human feedback, atomic checkpointing for crash recovery, and a single-call API for onboarding new models at runtime.  We release the full codebase, experiment scripts, and pre-trained artifacts to enable both immediate enterprise deployment and reproducible research on non-stationary LLM routing.